\section{Limitations and Ethics}

\paragraph{Scope.}
The benchmark supports claims about short-answer extractive QA. It should not
be used to claim that the same paradox, or the same policy ranking, holds for
long-form synthesis, agentic workflows, or domain-specific retrieval without
additional controls.

\paragraph{Faithfulness metrics.}
The multi-metric result is a strength and a limitation. It shows that evaluator
choice matters, but it does not establish which evaluator is correct. No
completed two-rater human calibration is reported; a 99-item template is staged
for future calibration.

\paragraph{Thresholds and CCS.}
Thresholds do not transfer cleanly across datasets: SQuAD, PubMedQA, and
NaturalQuestions show large in-domain versus off-diagonal recovery gaps. CCS is
also not causal in the matched-similarity test and has regimes where it is
uninformative about support, including coherent context sets that omit the
answer. It should be treated as one diagnostic among several.

\paragraph{Self-RAG harness.}
Self-RAG underperforms in our short-answer harness. This should be reported
carefully: the result is evidence about this matched inference setup, not a
global claim that Self-RAG is weak. Different prompting, decoding, or task
formats could change the outcome.

\paragraph{Compute and reproducibility.}
The revision was run under a zero-dollar constraint using local models and free
GPU sessions. That improves accessibility but leaves one planned independent
70B reproduction blocked. The earlier 70B observation is therefore not treated
as an independently rerun revision result. The artifact should disclose this
plainly rather than replace it with paid or unverifiable runs.

\paragraph{Ethical considerations.}
The benchmark is designed to reduce hallucination overclaims in RAG systems,
which is especially relevant for domains where unsupported answers can cause
harm. The main ethical risk is misuse: a single metric or a single dataset row
could be cherry-picked to market a RAG system as faithful. The benchmark
therefore emphasizes per-query logs, multiple metrics, and cost-aware reporting.
